\chapter{Eigenmodes: Resonant Frequencies and Q}
\label{ch:eigenmodes}

\chapterorient{We have studied electrostatics (\cref{ch:electrostatics}) as the
anchor pipeline: \texttt{assemble} an operator once, run a family of independent
solves over it, and \texttt{reduce} the results to a physical product. This
chapter takes that same skeleton and swaps exactly one corner---the
\texttt{solve}. Instead of a linear solve we write ourselves, the eigenmode
analysis hands the whole assembled problem to an opaque eigenvalue library and
receives the answer. It is the fourth solve corner from \cref{ch:situating}: the
black box. We will see that almost everything else about the pipeline is
unchanged, that the one Palace-written loop is the harmless post-processing
readout, and that the price of the black box is a subtlety---spurious
non-physical modes---that we must handle \emph{before} the solve, not after.}

\section{The physical question: how does a cavity ring?}

Strike a wine glass and it sings at a definite pitch. Pluck a guitar string and
it vibrates at a fundamental frequency plus a ladder of overtones. A closed
electromagnetic cavity---a hollow metal box, a superconducting resonator, a
dielectric puck---does exactly the same thing with electromagnetic fields. With
\emph{no external drive at all}, the cavity supports a discrete set of standing
electromagnetic oscillations, its \emph{resonant modes}. Each mode rings at its
own \emph{resonant frequency} $f$, decays at its own rate (captured by a
\emph{quality factor} $Q$), and has its own spatial pattern of $\vE$ and $\vB$
(\cref{fig:cavity-setup}). The eigenmode analysis computes the lowest few of
these---typically the handful of modes nearest some target frequency---because
those are the ones that matter for a filter passband, a qubit's readout
resonator, or an accelerator cavity.

\begin{physicsbox}
A \emph{closed} cavity is a region of space bounded (mostly) by conducting
walls, with no source inside and no port feeding energy in. Maxwell's equations
on such a domain are \emph{homogeneous}: there is no right-hand side. A
homogeneous linear system has only the trivial solution $\vE=\mathbf 0$
\emph{except} at special frequencies, where a nonzero field can sustain itself.
Those special frequencies are the resonant frequencies; the self-sustaining
fields are the modes. Physically, the cavity is a trap: a mode is a field
pattern that bounces between the walls and reinforces itself, neither growing nor
(in the ideal lossless case) decaying. Finding the modes is therefore not a
matter of \emph{driving} the cavity and seeing what comes out---it is a matter of
asking \emph{at what frequencies can the cavity oscillate on its own?}
\end{physicsbox}

\begin{figure}[t]
\centering
\begin{tikzpicture}[font=\footnotesize]
  % cavity walls
  \draw[thick, simInk, fill=simBgGray] (0,0) rectangle (6.2,3.0);
  \node[anchor=south west, font=\scriptsize, text=simGray] at (0.05,3.02)
    {conducting walls (PEC): $\vn\times\vE=\mathbf 0$};
  % standing-wave E field pattern (vertical arrows, sin envelope)
  \foreach \x in {0.4,0.8,...,5.8}{
    \pgfmathsetmacro{\amp}{1.05*sin(deg(pi*\x/6.2))}
    \draw[->, simBlue, thick] (\x,1.5) -- (\x,{1.5+\amp});
    \draw[->, simBlue, thick] (\x,1.5) -- (\x,{1.5-\amp});
  }
  % envelope
  \draw[densely dashed, simRust, semithick, domain=0:6.2, samples=80, smooth]
    plot (\x,{1.5+1.15*sin(deg(pi*\x/6.2))});
  \draw[densely dashed, simRust, semithick, domain=0:6.2, samples=80, smooth]
    plot (\x,{1.5-1.15*sin(deg(pi*\x/6.2))});
  \node[text=simBlue] at (3.1,2.85) {$\vE$};
  \node[text=simRust, anchor=west, font=\scriptsize] at (6.25,2.4)
    {standing-wave envelope};
  % half-wavelength bracket
  \draw[<->, simGreen, thick] (0,0.25) -- (6.2,0.25)
    node[midway, fill=simBgGray, inner sep=1pt, text=simGreen]
      {$L=\tfrac{1}{2}\lambda$ (fundamental)};
\end{tikzpicture}
\caption[A cavity's lowest mode]{The fundamental mode of a closed cavity. With no
external drive, the field cannot escape the conducting walls (the perfect-electric-conductor
boundary forces the tangential $\vE$ to vanish there, \cref{ch:bc}), so it forms a
standing wave: an electric-field pattern (blue) whose amplitude follows a fixed
spatial envelope (dashed) and oscillates in time at the resonant frequency. The
fundamental fits exactly half a wavelength across the cavity, $L=\tfrac12\lambda$;
higher modes fit more half-wavelengths and ring at higher frequencies. The
eigenmode analysis computes these patterns and their frequencies directly, from
the geometry alone.}
\label{fig:cavity-setup}
\end{figure}

\subsection{The same skeleton, one corner swapped}

In \cref{ch:situating} we reduced every Palace analysis to one skeleton,
\[
  \texttt{config} \to \texttt{assemble} \to \texttt{solve}
  \to \texttt{reduce} \to \texttt{product},
\]
and we placed electrostatics (\cref{ch:electrostatics}) at its center as the
universal anchor. Eigenmode analysis is the \emph{same machine} with a single
part swapped---but the part it swaps is the most consequential one. In
electrostatics the \texttt{solve} stage is a \emph{fixed-operator map}: assemble
the stiffness $\Kop$ once, then call our own linear solver $\Kop\vx_i=\vb_i$ once
per conductor, in a loop we write and control (\code{solve\_family}). In the
eigenmode analysis the \texttt{solve} stage is not a linear solve at all, and
there is no loop we write: we assemble the operator pencil, hand it whole to an
eigenvalue library, and get the converged modes back from a single call.

\begin{keyidea}
The eigenmode pipeline is the anchor pipeline with the solve corner replaced.
Electrostatics solves $\Kop\vx_i=\vb_i$ in a Palace-written loop; the eigenmode
analysis solves the \emph{eigenproblem} $\Kop\vx=\lambda\Mop\vx$, and the entire
iteration that does so lives \emph{inside an opaque library}. Palace writes no
solve loop here at all. \emph{The only loop Palace writes in the eigenmode driver
is the post-processing readout map over the modes the library already
converged.} This is the fourth solve corner of \cref{ch:situating}: the opaque
black box.
\end{keyidea}

The rest of this chapter walks the three stages in order---assemble, the
black-box solve, reduce---then handles the one subtlety the black box forces on
us (spurious modes), and closes by reading out the mode fields and naming a few
applications.

\section{Assemble: the eigenproblem pencil, built once}
\label{sec:eig-assemble}

The assemble stage is almost unchanged from the anchor. We build operators by
folding weak-form terms over the finite-element space (\cref{ch:fem}); the only
difference is \emph{which} operators, and that there is more than one.

The fields live in the Nédélec edge-element space $\Hcurl$ (\cref{ch:functionspaces}),
the natural home for $\vE$: it builds tangential continuity across element faces
into the basis, exactly matching the physics of an electric field at a material
interface. On that space we assemble up to three operators:
\begin{itemize}[nosep]
  \item the \textbf{curl--curl stiffness} $\Kop$, the discretization of
  $\Curl(\nu\,\Curl\,\cdot\,)$ where $\nu=1/\mu$ is the reluctivity---this is the
  ``spring constant'' of the electromagnetic field, the operator that resists
  curvature;
  \item the \textbf{mass} $\Mop$, the discretization of $\varepsilon\,(\cdot)$
  weighted by permittivity---the ``inertia'' of the field;
  \item optionally the \textbf{damping} $\Cop$, present only when the cavity is
  \emph{lossy} (finite wall conductivity, lossy dielectric, or a resistive
  port)---the operator that drains energy.
\end{itemize}
All three come from the same assemble fold we have used since \cref{ch:fem}; in
the source they are three calls,
\lstinline[style=l4style]|GetStiffnessMatrix|,
\lstinline[style=l4style]|GetDampingMatrix|, and
\lstinline[style=l4style]|GetMassMatrix|, each run \emph{once}.

\subsection{Lossless: the generalized eigenproblem}

When the cavity is lossless, $\Cop$ is absent and Maxwell's source-free
time-harmonic equation, discretized, becomes a \emph{generalized eigenvalue
problem} (\cref{ch:reductions-pde}):
\begin{equation}
  \Kop\,\vx \;=\; \lambda\,\Mop\,\vx,
  \qquad \lambda = \omega^2,
  \label{eq:gevp}
\end{equation}
where $\vx$ is the vector of degrees of freedom of the mode field $\vE$, the
eigenvalue $\lambda$ is the squared angular frequency $\omega^2$, and the
eigenvector $\vx$ \emph{is} the discretized mode. This is exactly the generalized
EVP we met abstractly in \cref{ch:reductions-pde} and studied numerically in
\cref{ch:eigenvalues}; here it arrives with a concrete physical reading. Each
eigenpair $(\lambda,\vx)$ is one resonance: $\omega=\sqrt\lambda$ is its angular
frequency, $\vx$ is its field pattern.

\begin{notationbox}
The eigenvalue is the \emph{square} of the angular frequency, $\lambda=\omega^2$,
not the frequency itself. This is the single most important bookkeeping fact in
the chapter. It means the solver returns $\lambda$, and the readout must
\emph{un-transform} it---$\omega=\sqrt\lambda$---to get a frequency, then divide
by $2\pi$ to get the ordinary frequency $f=\omega/2\pi$ in hertz. Forgetting the
square root, or the $2\pi$, is the classic first-run-off-by-a-lot bug. We meet
the un-transform again in \cref{sec:eig-reduce}.
\end{notationbox}

\subsection{Lossy: the quadratic eigenproblem}

When the cavity loses energy, the damping operator $\Cop$ enters and the problem
becomes \emph{quadratic} in the eigenvalue (\cref{ch:reductions-pde}):
\begin{equation}
  \bigl(\Kop + \lambda\,\Cop + \lambda^2\,\Mop\bigr)\,\vx \;=\; \mathbf 0,
  \qquad \lambda = i\,\omega.
  \label{eq:qevp}
\end{equation}
Now $\lambda$ is not $\omega^2$ but $i\omega$, so a lossy mode's eigenvalue is
\emph{complex}: its imaginary part carries the frequency and its real part the
decay rate. The quadratic form is what lets a mode decay---the term linear in
$\lambda$ couples frequency to loss. We do not solve \cref{eq:qevp} by hand;
\cref{ch:eigenvalues} showed how a quadratic eigenproblem is linearized into a
larger ordinary one that the same library swallows. For our purposes the
quadratic case differs from the linear case only in (i) which operators the
pencil carries, and (ii) how the eigenvalue un-transforms to $\omega$. Both
differences are absorbed into the single \emph{problem-type} parameter the solve
takes; the shape of the pipeline does not change.

\begin{definition}[The operator pencil]
The \emph{pencil} is the bundle of assembled operators handed to the eigensolver:
the pair $(\Kop,\Mop)$ for the linear generalized eigenproblem \cref{eq:gevp}, or
the triple $(\Kop,\Cop,\Mop)$ for the quadratic eigenproblem \cref{eq:qevp},
together with a spectral-transform target $\sigma$ (\cref{ch:eigenvalues}) telling
the solver \emph{which} part of the spectrum we want. The pencil is assembled
exactly once; nothing in the solve mutates it.
\end{definition}

\begin{keyidea}
Assembly produces an operator \emph{pencil}, not a single matrix: $(\Kop,\Mop)$
for a lossless cavity (the generalized EVP $\Kop\vx=\lambda\Mop\vx$,
$\lambda=\omega^2$) or $(\Kop,\Cop,\Mop)$ for a lossy one (the quadratic EVP
$(\Kop+\lambda\Cop+\lambda^2\Mop)\vx=\mathbf 0$, $\lambda=i\omega$). It is built
\emph{once}, exactly like the single $\Kop$ of electrostatics---the assemble
corner is shared; only the count and the meaning of the operators differ.
\end{keyidea}

\section{Solve: the opaque black box}
\label{sec:eig-solve}

Here is where the pipeline diverges sharply from the anchor. We have the pencil.
What do we do with it?

In electrostatics we would now write a loop: for each right-hand side, call our
own conjugate-gradient or direct solver, collect the answer. We control that
loop; we can see every linear solve it performs; the iteration is
\emph{ours}. The eigenmode analysis does nothing of the kind. We make a single
call,
\[
  \texttt{eigs} \;=\; \texttt{eigsolve}\;\bigl(\text{pencil},\ \sigma,\
  n_{\text{modes}}\bigr),
\]
handing the entire assembled problem---operators, target, requested mode
count---to a specialized eigenvalue library (SLEPc or ARPACK in Palace's case,
\cref{ch:eigenvalues}). The library runs its own sophisticated iteration, a
Krylov--Schur or Arnoldi process with a shift-and-invert spectral transform
(\cref{ch:krylovschur}), and returns the converged eigenpairs. We do not write
that iteration. We do not see it. We see one call and its result
(\cref{fig:eig-pipeline}).

\begin{figure}[t]
\centering
\begin{tikzpicture}[node distance=7mm and 10mm, font=\small]
  \node[data] (cfg) {config};
  \node[stage, right=of cfg, align=center] (asm)
    {assemble\\pencil};
  \node[extern, right=of asm, minimum width=22mm, minimum height=11mm,
        align=center] (solve) {\textbf{eigsolve}\\\scriptsize SLEPc / ARPACK};
  \node[stage, right=of solve, align=center] (map) {map\\readout};
  \node[product, right=of map, align=center] (prod) {$(f,Q)$\\table};
  \draw[flow] (cfg) -- (asm);
  \draw[flow] (asm) -- node[above, font=\scriptsize]{$(\Kop,\Cop,\Mop),\sigma$} (solve);
  \draw[flow] (solve) -- node[above, font=\scriptsize]{eigenpairs} (map);
  \draw[flow] (map) -- (prod);
  % the library's hidden inner loop: callback for operator apply
  \draw[refedge] (solve.south) ++(0,-1mm)
    .. controls +(0,-7mm) and +(0,-7mm) .. ($(solve.south)+(0,-1mm)$)
    node[midway, below, font=\scriptsize, text=simViolet, align=center]
      {hidden inner loop:\\Krylov--Schur + shift-invert\\(calls back for $\mathbf{A}\vx$)};
  % the one Palace-written loop annotation
  \node[font=\scriptsize, text=simGreen, align=center, below=2mm of map]
    {the \emph{only} loop\\Palace writes};
\end{tikzpicture}
\caption[The eigenmode pipeline]{The eigenmode pipeline, drawn to contrast with
the electrostatic anchor (\cref{ch:electrostatics}). Assembly builds the operator
pencil once. The \texttt{eigsolve} box is dashed and violet---the convention of
\cref{ch:situating} for a corner where \emph{we write no loop}: the entire
eigen-iteration (Krylov--Schur with shift-and-invert, \cref{ch:krylovschur})
happens inside the library, which only calls back out to us for the operator
application $\mathbf A\vx$. The single Palace-written loop is the post-processing
\texttt{map readout} (green) that turns each converged eigenpair into a row of the
$(f,Q)$ table. Compare \cref{ch:electrostatics}, where the solve box is a solid
Palace-written \code{solve\_family} loop instead of a black box.}
\label{fig:eig-pipeline}
\end{figure}

\subsection{The black box still calls back}

``Opaque'' does not mean ``self-contained.'' The library does not know what
$\Kop$, $\Cop$, or $\Mop$ \emph{are}; it knows only that it can ask us to apply
them. The shift-and-invert spectral transform (\cref{ch:eigenvalues}) needs to
apply the operator $(\Kop-\sigma\Mop)^{-1}\Mop$ to a vector on every inner step,
and \emph{that inverse is itself a linear solve}. So the library, mid-iteration,
calls back into our code asking us to apply the operator---and to apply that
inverse we run one of our own linear solves (the very \code{ksp\_solve}
machinery from electrostatics, \cref{ch:electrostatics}). The picture is a black
box with a wire poking out: the library owns the outer eigen-iteration, but the
operator applications---including an inner linear solve buried inside the
shift-and-invert---are still ours. We see those applications; we do not see how
the library schedules them or decides it has converged.

\begin{keyidea}
The eigensolver is opaque but not isolated: it runs the outer eigen-iteration
itself, and \emph{calls back} into our code for each operator application
$\mathbf A\vx$. Because the spectral transform applies $(\Kop-\sigma\Mop)^{-1}$,
each callback contains an \emph{inner linear solve}---the same linear-solve
machinery electrostatics uses, now nested one level down, inside someone else's
loop. The opacity is precise: we own the operator applies, the library owns the
iteration and the convergence test.

\end{keyidea}

\subsection{Why hand it off at all?}

It is fair to ask why we surrender control here when we kept it for every other
solve corner. The answer is honest engineering, not laziness. A robust
eigenvalue iteration that finds a few interior eigenpairs of a large
non-symmetric (in the lossy case) pencil, restarts to control memory, locks
converged pairs, and deflates against them is a substantial piece of numerical
software---Krylov--Schur with all its restart and purging logic
(\cref{ch:krylovschur}). Palace's authors made the deliberate choice to call a
mature library (SLEPc, itself wrapping PETSc; or ARPACK) rather than rewrite it.
Our framework records that choice faithfully: the \code{eigsolve} combinator is
\emph{deliberately} opaque. It names the library boundary, states the contract
across it (in: a pencil and a target; out: converged eigenpairs and a count), and
stops.

\begin{pitfall}
The opacity of \code{eigsolve} is a property of \emph{Palace's} code, not of the
mathematics. It is tempting to read ``black box'' as ``we do not understand the
eigen-iteration.'' We do---\cref{ch:eigenvalues} and \cref{ch:krylovschur} build
Krylov--Schur from scratch. The black box marks where \emph{Palace} stops writing
its own loop and calls out, not where understanding stops. A reader who wants to
see inside the box opens those two chapters, which reconstruct exactly what the
library does, in our own vocabulary. The honest record is: \emph{Palace} treats
it as a black box; the book does not.
\end{pitfall}

\subsection{Contrast with the anchor's solve}

The contrast with electrostatics is worth drawing as a side-by-side
(\cref{fig:solve-vs-map}), because it is the whole point of the chapter. In
electrostatics the solve stage is a loop \emph{we} write, the
\code{solve\_family} fixed-operator map: visible, controllable, one linear solve
per right-hand side. In the eigenmode analysis the solve stage is a single
library call, and the only loop we write comes \emph{after} the solve---the
readout map over the modes the library has already produced. The difference is
not which equation we solve; it is who owns the iteration.

\begin{figure}[t]
\centering
\begin{tikzpicture}[font=\footnotesize]
  % --- left: electrostatic solve (a loop we write) ---
  \begin{scope}[shift={(0,0)}]
    \node[font=\small\bfseries, align=center] at (1.6,2.55)
      {Electrostatic: a loop \emph{we} write};
    \node[opbox, minimum width=8mm] (k) at (0.1,1.2) {$\Kop$};
    \foreach \i/\y in {1/1.85,2/1.2,3/0.55}{
      \node[product, minimum width=5mm, minimum height=4mm, scale=0.7]
        (es\i) at (2.1,\y) {};
      \draw[flow, shorten >=1pt] (k.east) -- (es\i.west);}
    \node[font=\scriptsize\ttfamily, text=simGreen] at (1.6,-0.05)
      {solve\_family (corner 1)};
    \node[font=\scriptsize\itshape, align=center] at (1.6,-0.55)
      {one linear solve per RHS,\\every solve visible};
  \end{scope}
  % divider
  \draw[simGray, densely dashed] (4.0,-0.9) -- (4.0,2.8);
  % --- right: eigenmode solve (a black box) ---
  \begin{scope}[shift={(4.6,0)}]
    \node[font=\small\bfseries, align=center] at (1.6,2.55)
      {Eigenmode: a black box};
    \node[opbox, minimum width=8mm, align=center] (pen) at (0.0,1.2)
      {$(\Kop,$\\$\Cop,\Mop)$};
    \node[extern, minimum width=13mm, minimum height=9mm] (bb) at (1.9,1.2)
      {\scriptsize eigsolve};
    \draw[flow, shorten >=1pt] (pen.east) -- (bb.west);
    \node[product, minimum width=8mm, minimum height=5mm, scale=0.75]
      (out) at (3.35,1.2) {\scriptsize eigs};
    \draw[flow, shorten >=1pt] (bb.east) -- (out.west);
    \node[font=\scriptsize\ttfamily, text=simViolet] at (1.6,-0.05)
      {eigsolve (corner 4)};
    \node[font=\scriptsize\itshape, align=center] at (1.6,-0.55)
      {one call, iteration hidden;\\our loop is the readout, after};
  \end{scope}
\end{tikzpicture}
\caption[Solve corner vs.\ map corner]{The two solve stages side by side. Left:
electrostatics runs a \code{solve\_family} map---a Palace-written loop of
independent linear solves over the fixed operator $\Kop$, every solve visible
(\cref{ch:electrostatics}). Right: the eigenmode analysis runs \code{eigsolve}---a
single call into an opaque library (dashed violet), whose internal iteration we do
not see. The eigenmode driver's \emph{own} loop is not in the solve at all; it is
the post-processing readout map that follows. This is the structural heart of the
chapter: same skeleton, the solve corner swapped from a loop-we-write to a
loop-we-do-not.}
\label{fig:solve-vs-map}
\end{figure}

\section{Keeping the modes physical: the gradient null-space}
\label{sec:eig-spurious}

There is a subtlety that the black box forces us to confront, and it is not
optional. The curl--curl stiffness operator $\Kop$ has a large \emph{null
space}: for \emph{any} scalar potential $\phi$, the gradient field $\Grad\phi$
satisfies $\Curl\Grad\phi=\mathbf 0$ identically, so $\Kop(\Grad\phi)=\mathbf 0$.
On the discrete space this means $\Kop$ has a whole subspace of eigenvectors with
eigenvalue \emph{zero}---one for every discrete gradient. These are not
resonances. They are an artifact of the curl--curl operator's structure, the
discrete shadow of the Helmholtz decomposition (\cref{ch:bc}): every $\Hcurl$
field splits into a curl part (the physics) and a gradient part (the null space),
and only the curl part carries a real mode.

If we hand the raw pencil to the eigensolver and ask for the modes nearest zero
frequency, it will dutifully return a cloud of these spurious zero and
near-zero ``modes''---numerically polluted gradients---and bury the real
low-frequency resonances among them (\cref{fig:spurious}). The black box cannot
tell physics from artifact; it returns eigenpairs, and these are genuine
eigenpairs of the matrix we gave it. The fix has to happen on \emph{our} side,
\emph{before} the solve.

\begin{figure}[t]
\centering
\begin{tikzpicture}[font=\footnotesize]
  % --- before: spectrum with a spurious cloud near 0 ---
  \begin{scope}[shift={(0,0)}]
    \node[font=\small\bfseries] at (3.0,1.9) {Before projection};
    \draw[->, thick, simInk] (0,0) -- (6.2,0) node[right]{$\omega$};
    \draw[thick, simInk] (0,-0.1) -- (0,0.1) node[above, font=\scriptsize]{$0$};
    % spurious cloud near 0
    \foreach \x in {0.05,0.12,0.2,0.28,0.35,0.45,0.55,0.62,0.7,0.8}{
      \fill[simRust] (\x,0) circle (1.6pt);}
    \node[text=simRust, font=\scriptsize, align=center] at (0.5,0.75)
      {spurious gradient\\null-space modes};
    \draw[simRust, ->] (0.5,0.5) -- (0.4,0.12);
    % real modes
    \foreach \x in {2.4,3.6,4.9}{
      \fill[simBlue] (\x,0) circle (2.4pt);}
    \node[text=simBlue, font=\scriptsize] at (3.6,0.55) {real resonances};
    \node[text=simRust, font=\scriptsize, align=center] at (3.0,-0.7)
      {the real modes are \emph{buried} in the cloud};
  \end{scope}
  % --- after: clean spectrum ---
  \begin{scope}[shift={(0,-2.6)}]
    \node[font=\small\bfseries] at (3.0,1.9) {After divergence-free projection};
    \draw[->, thick, simInk] (0,0) -- (6.2,0) node[right]{$\omega$};
    \draw[thick, simInk] (0,-0.1) -- (0,0.1) node[above, font=\scriptsize]{$0$};
    % only real modes
    \foreach \x in {2.4,3.6,4.9}{
      \fill[simBlue] (\x,0) circle (2.4pt);}
    \node[text=simBlue, font=\scriptsize] at (3.6,0.55) {real resonances};
    \node[text=simGreen, font=\scriptsize, align=center] at (0.5,0.7)
      {null space\\removed};
    \draw[simGreen, ->] (0.5,0.42) -- (0.18,0.1);
    \node[text=simGreen, font=\scriptsize, align=center] at (3.0,-0.7)
      {the solver now returns physical modes first};
  \end{scope}
\end{tikzpicture}
\caption[Spurious modes and the projection]{The gradient null-space and its
removal. Top: the curl--curl operator $\Kop$ annihilates every gradient field
($\Kop\Grad\phi=\mathbf 0$), so its raw spectrum carries a cloud of spurious
zero/near-zero ``modes'' (rust) that are numerical gradients, not resonances; the
real low-frequency resonances (blue) are buried among them, and the solver,
asked for the modes nearest zero, returns the artifacts first. Bottom: the
\emph{divergence-free projection} (\cref{ch:bc}) restricts the solve to the
gradient-free subspace, deleting the null-space cloud so the eigensolver returns
the physical resonances directly. The projection is wired in \emph{before} the
solve, because the black box cannot tell artifact from physics on its own.}
\label{fig:spurious}
\end{figure}

The fix is the \emph{divergence-free projection} of \cref{ch:bc}. We restrict the
solve to the subspace of fields orthogonal to all gradients---the divergence-free
(curl) part of the Helmholtz decomposition---so the null-space directions are
simply unavailable to the eigensolver. Concretely, Palace builds a projector $P$
onto the divergence-free subspace and wires it into the eigensolver: the starting
vector is projected, and the projector is applied so the iteration never strays
into the gradient subspace. With the null space projected out, the spectrum the
library sees has no spurious cloud, and the modes it returns are physical
resonances.

\begin{keyidea}
The curl--curl operator's gradient null space ($\Kop\Grad\phi=\mathbf 0$ for
every scalar $\phi$, \cref{ch:bc}) produces a cloud of spurious zero-frequency
``modes'' that are not resonances. The \emph{divergence-free projection} restricts
the solve to the gradient-free subspace, removing them, so the black box returns
physical modes. This projection is wired in \emph{before} the opaque solve---it
must be, because the black box cannot distinguish artifact from physics once it
is running. Handling the null space is the one place where the eigenmode pipeline
demands real physics knowledge the anchor never needed.
\end{keyidea}

\section{Reduce: the eigenfrequency-and-Q table}
\label{sec:eig-reduce}

The library returns converged eigenpairs $(\lambda_m,\vx_m)$. The reduce stage
turns them into the physical product the user asked for: a small \emph{table},
one row per mode (\cref{ch:reductions}, \cref{ch:situating})---the rank-1
``Shape~3'' reduction. Each row carries the mode's resonant frequency $f_m$ and
its quality factor $Q_m$. This stage \emph{is} the one Palace-written loop in the
driver: a \code{map} over the converged modes, with no coupling between rows
(mode~3's row does not depend on mode~2's), which is exactly what makes it a map
and not a fold.

\subsection{From eigenvalue to frequency: the un-transform}

The first column of the table is the resonant frequency, and getting it right is
pure bookkeeping---the un-transform of the eigenvalue back to a physical
frequency (\cref{ch:eigenvalues}). Two things must be undone. First, the
\emph{problem-type} transform from \cref{sec:eig-assemble}: for the linear EVP
the eigenvalue is $\lambda=\omega^2$, so $\omega=\sqrt\lambda$; for the quadratic
EVP it is $\lambda=i\omega$, so $\omega=\lambda/i$. Second, any \emph{spectral
shift} $\sigma$ used to target the interior of the spectrum
(\cref{ch:eigenvalues}) must be inverted---the solver works in shifted
coordinates, and the physical eigenvalue is recovered by undoing the shift,
exactly as in the shift-and-invert worked example of \cref{ch:eigenvalues}. With
$\omega$ in hand, the eigenfrequency is its real part over $2\pi$:
\begin{equation}
  f_m \;=\; \frac{\Re\,\omega_m}{2\pi}.
  \label{eq:freadout}
\end{equation}

\subsection{From the complex eigenvalue to Q}

The second column is the quality factor, and it lives entirely in the imaginary
part. A lossless cavity's $\omega$ is purely real and the mode rings forever; a
lossy cavity's $\omega$ has a positive imaginary part, the decay rate, and the
mode rings down. The quality factor measures how many radians of oscillation the
mode completes per e-folding of its amplitude. In the energy/loss convention
(\cref{ch:reductions}),
\begin{equation}
  Q_m \;=\; \frac{\omega_m}{\kappa_m}
        \;=\; \frac{\Re\,\omega_m}{2\,\Im\,\omega_m},
  \label{eq:qreadout}
\end{equation}
where $\kappa_m$ is the mode's loss rate. The factor of $2$ converts the
amplitude decay rate ($\Im\,\omega$) to the energy decay rate (energy is
amplitude squared, so it decays twice as fast); different texts place the $2$
differently, which is precisely why the reduction takes $\kappa$ as a parameter
rather than hard-coding the convention.

\begin{notationbox}
The loss rate $\kappa_m$ is not always read straight off $\Im\,\omega$. When the
loss enters through a specific physical channel---a resistive lumped port, say---Palace
computes $\kappa_m$ as a \emph{participation ratio}: the fraction of the mode's
energy dissipated in that channel, $\kappa_m=\tfrac12 R\,\abs{I_m}^2/E_m$ (resistor
self-energy over total mode energy), which then enters \cref{eq:qreadout} as
$Q_m=\omega_m/\kappa_m$. The two readings agree---$\Im\,\omega$ \emph{is} the
total loss rate---but the participation form attributes the loss to its source,
which is what a designer wants to know. Either way $Q$ is a scalar ratio, one
number per mode.
\end{notationbox}

\subsection{The lossless limit}

The honest edge case is the lossless cavity: $\Im\,\omega\to 0$, $\kappa\to 0$,
and \cref{eq:qreadout} divides by zero. The physical answer is $Q\to\infty$: an
ideal lossless resonator rings forever. The reduction does not produce a
\texttt{NaN}; it tests $\kappa=0$ and returns $Q=\infty$ directly, the one piece
of control flow in an otherwise purely arithmetic readout. The infinite-$Q$ case
is real, not exceptional: it is the limit every low-loss design---a
superconducting cavity, a high-$Q$ dielectric resonator---approaches, and the
table reports it faithfully.

\begin{keyidea}
The reduce stage is a \code{map} over the converged eigenpairs producing a
rank-1 table, one $(f,Q)$ row per mode. The frequency is the eigenvalue
\emph{un-transformed} ($\omega=\sqrt\lambda$ linear, $\lambda/i$ quadratic; shift
undone) and read as $f=\Re\,\omega/2\pi$. The quality factor is the
energy/loss ratio $Q=\Re\,\omega/(2\,\Im\,\omega)=\omega/\kappa$, with the
$\kappa\to0$ lossless limit handled as $Q=\infty$ rather than a division error.
This is the \emph{only} loop the eigenmode driver writes, and it is a pure
post-processing map---never a solve.
\end{keyidea}

\section{Worked examples}

\begin{workedexample}{The modes of a 1-D cavity, by hand}{cavity1d}
We compute the lowest resonances of the simplest cavity: a one-dimensional
segment of length $L$ with conducting ends, the electromagnetic analog of a
string fixed at both ends. The continuous problem is
\[
  -u''(x) = \lambda\, u(x), \qquad u(0)=u(L)=0, \qquad \lambda = \omega^2/c^2,
\]
whose exact modes are $u_n(x)=\sin(n\pi x/L)$ with $\lambda_n=(n\pi/L)^2$, giving
the familiar resonant-frequency ladder
\[
  f_n \;=\; \frac{c}{2\pi}\sqrt{\lambda_n} \;=\; \frac{n\,c}{2L},
  \qquad n=1,2,3,\dots
\]
The fundamental fits half a wavelength across the cavity ($f_1=c/2L$,
\cref{fig:cavity-setup}); each higher mode adds one more half-wavelength. We will
recover this ladder as the eigenvalues of a \emph{discrete} pencil
$\Kop\vx=\lambda\Mop\vx$, exactly the matrices the eigensolver would receive.

\emph{Discretize.} Put $N$ interior nodes at spacing $h=L/(N+1)$ and use the
standard second-difference stencil for $-u''$ and the (lumped) identity for the
mass. With $N=3$ ($h=L/4$), the stiffness and mass matrices are
\[
  \Kop = \frac{1}{h^2}\!\begin{pmatrix} 2 & -1 & 0\\ -1 & 2 & -1\\ 0 & -1 & 2\end{pmatrix},
  \qquad
  \Mop = \mathsf{I}_3,
\]
so the pencil reduces to the ordinary eigenproblem $\Kop\vx=\lambda\vx$ for the
tridiagonal $\Kop$.

\emph{Solve the $3\times3$ eigenproblem.} The eigenvalues of the symmetric
tridiagonal $\begin{psmallmatrix}2&-1&0\\-1&2&-1\\0&-1&2\end{psmallmatrix}$ are
known in closed form, $2-2\cos\!\bigl(\tfrac{k\pi}{4}\bigr)$ for $k=1,2,3$:
\[
  2-\sqrt2 \approx 0.586,\qquad 2,\qquad 2+\sqrt2 \approx 3.414.
\]
Dividing by $h^2=(L/4)^2=L^2/16$ gives the discrete $\lambda_k$:
\[
  \lambda_1 \approx \frac{9.37}{L^2},\quad
  \lambda_2 = \frac{32}{L^2},\quad
  \lambda_3 \approx \frac{54.6}{L^2}.
\]
\emph{Compare to the exact ladder.} The exact values are
$\lambda_n=(n\pi/L)^2 = \{9.87, 39.5, 88.8\}/L^2$. The lowest discrete eigenvalue,
$9.37/L^2$, is within $5\%$ of the exact $9.87/L^2$---a coarse three-node mesh
already pins the fundamental. The higher discrete eigenvalues sag below the exact
ones (the familiar dispersion error of a coarse mesh: short-wavelength modes are
underresolved), and they would tighten toward the ladder as $N$ grows. The
\emph{eigenvectors} of the $3\times3$ matrix are
$\vx_k=(\sin\tfrac{k\pi}{4},\sin\tfrac{2k\pi}{4},\sin\tfrac{3k\pi}{4})$---the
samples of $\sin(k\pi x/L)$ at the interior nodes, the discrete standing waves.

\emph{Read out the frequency.} Take the fundamental. Un-transform
$\lambda_1$ to a frequency: $\omega_1=c\sqrt{\lambda_1}=c\sqrt{9.37}/L=3.06\,c/L$,
so $f_1=\omega_1/2\pi=0.487\,c/L$, against the exact $f_1=c/2L=0.500\,c/L$. The
discrete eigenvalue, un-transformed exactly as \cref{sec:eig-reduce} prescribes,
lands within $3\%$ of the physical resonant frequency---and \emph{this little
$3\times3$ problem is, in miniature, precisely what the black box does at scale}:
build $(\Kop,\Mop)$, find $\lambda$, un-transform to $f$. The only difference at
full size is the matrix dimension and the fact that we hand the eigenproblem to a
library instead of reading the eigenvalues off a known formula.
\end{workedexample}

\begin{workedexample}{The Q of a lossy mode, and the lossless limit}{qfactor}
A lossy eigenmode solve returns a converged eigenpair whose recovered angular
frequency, after the un-transform of \cref{sec:eig-reduce}, is the complex number
\[
  \omega_m \;=\; \omega_r + i\,\omega_i
            \;=\; 2\pi\,(\SI{6.000}{\giga\hertz}) \;+\; i\,\bigl(2\pi\cdot\SI{1.5}{\mega\hertz}\bigr).
\]
The real part is a \SI{6}{\giga\hertz} oscillation; the imaginary part is the
amplitude decay rate. We want the table row $(f_m, Q_m)$.

\emph{Frequency.} By \cref{eq:freadout}, the frequency is the real part over
$2\pi$:
\[
  f_m \;=\; \frac{\Re\,\omega_m}{2\pi} \;=\; \SI{6.000}{\giga\hertz}.
\]

\emph{Quality factor.} By \cref{eq:qreadout} with the energy/loss convention
$Q=\Re\,\omega/(2\,\Im\,\omega)$,
\[
  Q_m \;=\; \frac{\omega_r}{2\,\omega_i}
        \;=\; \frac{2\pi\cdot 6.000\times 10^{9}}{2\cdot 2\pi\cdot 1.5\times 10^{6}}
        \;=\; \frac{6.000\times 10^{9}}{3.0\times 10^{6}}
        \;=\; 2000 .
\]
A $Q$ of $2000$ is a sharp resonance: the mode completes about $2000$
radians---roughly $320$ full cycles---before its amplitude falls by a factor of
$e$. The factor of $2$ in the denominator is the amplitude-to-energy conversion
(\cref{sec:eig-reduce}); had we used the bare $\omega_r/\omega_i$ convention we
would have reported $4000$, which is why the reduction carries $\kappa$ as a
parameter and never hard-codes the factor.

\emph{The lossless limit.} Now let the loss vanish, $\omega_i\to 0$: a perfect
cavity. The denominator of \cref{eq:qreadout} goes to zero and $Q_m\to\infty$.
The reduction does not divide by zero and emit a \texttt{NaN}; it tests
$\kappa=0$ and returns $Q=\infty$, the honest physical answer---a lossless cavity
rings forever. This single guard is the one branch in an otherwise straight-line
arithmetic readout, and it exists because the infinite-$Q$ case is the ideal that
every superconducting or low-loss design approaches, not an error to be trapped.
\end{workedexample}

\section{Reading out the mode fields, and where this goes}
\label{sec:eig-fields}

The $(f,Q)$ table is the headline product, but each converged eigenpair also
carries the mode \emph{field}, and the readout map recovers it too. The
eigenvector $\vx_m$ \emph{is} the electric mode field $\vE_m$, in the
$\Hcurl$ degrees of freedom; the magnetic field follows from Faraday's law as the
curl of $\vE$,
\begin{equation}
  \vB_m \;=\; -\frac{1}{i\,\omega_m}\,\Curl\,\vE_m,
  \label{eq:bfield}
\end{equation}
one curl and a complex scale per mode. (An eigenvector is defined only up to a
complex scalar, so Palace fixes the phase by a normalization before recording
the field---otherwise the same mode could come back rotated arbitrarily in the
complex plane.) These fields are what a designer plots to \emph{see} the mode:
where the electric field piles up (the high-field regions that set the breakdown
limit), where the magnetic field loops, which part of the structure a given
resonance lives in.

The whole pipeline, finally, is the clean composition the framework names
(\cref{lst:eigcompose}): assemble the pencil, one black-box solve, one readout
map.

\begin{figure}[t]
\begin{lstlisting}[style=l4style, caption={[The eigenmode composition]The
eigenmode feature as a composition of firm vocabulary. Assemble the
$(\Kop,\Cop,\Mop)$ pencil once (three folds), hand it to the opaque
\texttt{eigsolve} black box once, and \texttt{map} the per-mode readout over the
converged eigenpairs. The only Palace-written loop is the final \texttt{map};
the solve is a single opaque call.}, label={lst:eigcompose}]
eigenmode :: EigenmodeConfig -> EigenmodeResult
eigenmode cfg =
  let space  = nd_space cfg                          -- H(curl) edge space
      k      = fe_assemble space [ curl_curl  (reluctivity cfg) ]   -- K, once
      c      = fe_assemble space [ damping    (conductivity cfg) ]  -- C, once (may be empty)
      m      = fe_assemble space [ mass_term  (permittivity cfg) ]  -- M, once
      pencil = eig_pencil k c m (target cfg) (n_modes cfg)          -- bundle + target
      eigs   = eigsolve pencil (initial_space cfg)   -- ONE opaque black-box solve
  in  map (readout cfg) eigs                          -- the only loop we write: readout map
\end{lstlisting}
\end{figure}

\begin{remark}[Applications]
The eigenmode analysis is the workhorse behind a surprising range of devices.
Microwave \emph{cavity filters} are built by coupling several resonators, each a
controlled mode whose $f$ and $Q$ this analysis predicts before any metal is cut.
Superconducting-qubit \emph{readout resonators} and \emph{coupling buses} are
high-$Q$ modes of on-chip cavities; the eigenmode solve gives their frequencies
and participation ratios (\cref{sec:eig-reduce}), which set qubit--cavity
coupling and dictate where loss must be controlled. Particle-accelerator
\emph{RF cavities} are designed mode-by-mode for a target accelerating frequency.
In every case the deliverable is the same rank-1 table: a short list of
$(f,Q)$ rows, plus the field patterns that explain them.
\end{remark}

A close cousin of this analysis appears one chapter on. The \emph{boundary-mode}
(wave-port) analysis of \cref{ch:waveports} is, in the taxonomy of
\cref{ch:situating}, eigenmode on a two-dimensional cross-section: the same
opaque black-box solve corner, the same rank-1 table shape, applied to the
transverse modes of a waveguide cross-section rather than the volume modes of a
cavity. Having understood the volume eigenmode, the wave-port analysis will read
as a small, well-localized change of domain.

\summarysection
The eigenmode analysis is the anchor pipeline (\cref{ch:electrostatics}) with the
\texttt{solve} corner swapped for the fourth corner of \cref{ch:situating}: the
opaque black box. \emph{Assemble} builds an operator \emph{pencil} once---$(\Kop,\Mop)$
for a lossless cavity (the generalized EVP $\Kop\vx=\lambda\Mop\vx$,
$\lambda=\omega^2$) or $(\Kop,\Cop,\Mop)$ for a lossy one (the quadratic EVP,
$\lambda=i\omega$). \emph{Solve} is a single \code{eigsolve} call into a
specialized library (SLEPc/ARPACK, Krylov--Schur with shift-and-invert,
\cref{ch:eigenvalues}, \cref{ch:krylovschur}); Palace writes no eigen-iteration
loop, and the library only calls back for the operator applies (each containing
an inner linear solve from the shift-invert). Before the solve, the
\emph{divergence-free projection} (\cref{ch:bc}) removes the curl--curl operator's
spurious gradient null-space ``modes,'' so the black box returns physical
resonances. \emph{Reduce} is a \code{map} over the converged eigenpairs producing
the rank-1 $(f,Q)$ table: $f=\Re\,\omega/2\pi$ after the eigenvalue un-transform
(\cref{ch:eigenvalues}, \cref{ch:reductions}), $Q=\Re\,\omega/(2\,\Im\,\omega)$
with the lossless $\kappa\to0$ limit reported as $Q=\infty$. The mode fields read
out as $\vE=\vx$ and $\vB=-\tfrac{1}{i\omega}\Curl\vE$. The single Palace-written
loop in the whole driver is that final readout map---never a solve.

\furtherreading
Pozar's \emph{Microwave Engineering}~\citep{pozar2011} is the standard
device-level account of cavity resonators, resonant frequency, and the quality
factor, including the energy/loss definition of $Q$ and the loaded-versus-unloaded
distinction we glossed over. Jin's \emph{Finite Element Method in
Electromagnetics}~\citep{jin2014} develops the curl--curl eigenproblem on
$\Hcurl$ edge elements in full, including the gradient null-space and the
divergence-free constraint that tames it (\cref{sec:eig-spurious})---the FEM-level
companion to this chapter. For the eigensolver itself---why a few interior
eigenpairs of a large pencil are best found by Krylov--Schur with a
shift-and-invert spectral transform, and why that is the right thing to hand to a
library---Saad's \emph{Numerical Methods for Large Eigenvalue
Problems}~\citep{saad2011eig} is the reference; \cref{ch:eigenvalues} and
\cref{ch:krylovschur} build the relevant machinery on its foundation.

\exercisessection
\begin{enumerate}[nosep]
  \item \textbf{The swapped corner.} State precisely which stage of the
  \texttt{config}$\to$\texttt{assemble}$\to$\texttt{solve}$\to$\texttt{reduce}
  skeleton differs between electrostatics (\cref{ch:electrostatics}) and the
  eigenmode analysis, and which stages are unchanged. In one sentence, name the
  single structural property of the solve stage that makes the eigenmode corner
  ``opaque.''
  \item \textbf{The only loop.} The chapter insists the eigenmode driver writes
  exactly one loop, and that it is not a solve. Which loop is it, what does each
  of its iterations produce, and why is it a \code{map} (independent rows) rather
  than a \code{fold} (chained state)?
  \item \textbf{Un-transform.} A linear-EVP solve returns the eigenvalue
  $\lambda=39.5/L^2$ (in units where $c=1$). Recover the ordinary frequency $f$.
  Then a quadratic-EVP solve returns $\lambda=i\,(2\pi\cdot 4.0\times10^9)$;
  recover $f$ from \emph{it}. Why do the two cases use different un-transforms?
  \item \textbf{Q and the factor of two.} A mode has
  $\omega=2\pi(\SI{10}{\giga\hertz})+i\,2\pi(\SI{5}{\mega\hertz})$. Compute $Q$ in
  the energy/loss convention $Q=\Re\,\omega/(2\,\Im\,\omega)$. A colleague reports
  $Q=2000$ for the same mode; what convention did they use, and which factor did
  they drop?
  \item \textbf{The lossless limit.} Explain why the reduce stage tests
  $\kappa=0$ and returns $Q=\infty$ rather than letting \cref{eq:qreadout} divide
  by zero. Is $Q=\infty$ a numerical degeneracy to be avoided, or a physical
  answer to be reported? Give one real device whose $Q$ this limit describes.
  \item \textbf{Spurious modes.} The curl--curl operator satisfies
  $\Kop\Grad\phi=\mathbf 0$ for every scalar field $\phi$. Explain why this gives
  the raw eigenproblem a cloud of zero-eigenvalue ``modes,'' why those are not
  resonances, and why the divergence-free projection must be applied
  \emph{before} the black-box solve rather than filtered out of its output
  afterward.
  \item \textbf{The callback.} The eigensolver is opaque, yet the chapter says it
  ``calls back'' into Palace's code mid-iteration. What does it call back to
  compute, and why does that callback contain an \emph{inner} linear solve?
  Relate the inner solve to the shift-and-invert spectral transform of
  \cref{ch:eigenvalues}.
  \item \textbf{Refining the mesh.} In \cref{wex:cavity1d} the coarse $3\times3$
  mesh gave $\lambda_1\approx 9.37/L^2$ against the exact $9.87/L^2$. Without
  recomputing, predict qualitatively what happens to the three discrete
  eigenvalues as $N$ grows from $3$ to $30$: which converge fastest, and why does
  the highest mode lag? (Hint: how many nodes per wavelength does each mode get?)
\end{enumerate}
